\section{Technical \#1: Exact}\label{sec:technical-1:-exact}

\begin{figure}
  \centering
  \resizebox{\columnwidth}{!}{%
  \begin{tikzpicture}[
      font=\small\sffamily,
      tbl/.style={draw, rounded corners=2pt, align=center, minimum width=1.9cm,
                  minimum height=0.7cm, font=\footnotesize\sffamily},
      stp/.style={draw, rounded corners=2pt, align=center, fill=black!4,
                  minimum width=2.6cm, minimum height=1.1cm, font=\scriptsize\sffamily},
      cl/.style={draw, rounded corners=2pt, align=center, minimum width=1.9cm,
                 minimum height=0.8cm, font=\scriptsize\sffamily},
      jn/.style={draw, rounded corners=2pt, align=center, fill=black!8,
                 minimum width=1.9cm, minimum height=0.8cm, font=\scriptsize\sffamily},
      red/.style={draw, rounded corners=2pt, align=center, minimum width=5.4cm,
                  minimum height=0.9cm, font=\scriptsize\sffamily},
      outp/.style={draw, rounded corners=2pt, align=center, minimum width=1.9cm,
                   minimum height=0.6cm, font=\footnotesize\sffamily},
      flow/.style={-{Latex[length=2mm]}, thick},
      setup/.style={-{Latex[length=1.6mm]}, black!45},
      grp/.style={font=\scriptsize\bfseries\sffamily, rotate=90},
    ]

    % ---- inputs ----
    \node[tbl] (corpus) at (1.5,10.2) {corpus};
    \node[tbl] (probe)  at (4.5,10.2) {probe};

    % ---- steps 1 & 2 ----
    \node[stp] (s1) at (1.5,8.4) {\textbf{Step 1}\\ cluster corpus\\ k-means into $N$ clusters};
    \node[stp] (s2) at (4.5,8.4) {\textbf{Step 2}\\ cocluster probe\\ each row to $p$ nearest clusters};

    % ---- clusters ----
    \node[cl] (c1) at (0,6.4) {cluster 1\\ \scriptsize corpus + probe};
    \node[cl] (c2) at (3,6.4) {cluster 2\\ \scriptsize corpus + probe};
    \node[cl] (c3) at (6,6.4) {cluster 3\\ \scriptsize corpus + probe};

    % ---- step 3: exact join, per cluster ----
    \node[jn] (j1) at (0,4.5) {\textbf{Step 3}\\ (in-cluster) exact join};
    \node[jn] (j2) at (3,4.5) {\textbf{Step 3}\\ (in-cluster) exact join};
    \node[jn] (j3) at (6,4.5) {\textbf{Step 3}\\ (in-cluster) exact join};

    % ---- step 4: reduction (local + global) ----
    \node[red] (r) at (3,2.6)
      {\textbf{Step 4}\\ local/global reduction};

    % ---- output ----
    \node[outp] (out) at (3,1.1) {result};

    % ---- setup edges (clustering fan) ----
    \draw[setup] (corpus) -- (s1);
    \draw[setup] (probe)  -- (s2);
    \foreach \c in {c1,c2,c3} {
      \draw[setup] (s1) -- (\c);
      \draw[setup] (s2) -- (\c);
    }

    % ---- main flow ----
    \draw[flow] (c1) -- (j1);
    \draw[flow] (c2) -- (j2);
    \draw[flow] (c3) -- (j3);
    \draw[flow] (j1) -- (r);
    \draw[flow] (j2) -- (r);
    \draw[flow] (j3) -- (r);
    \draw[flow] (r)  -- (out);

    % ---- grouping boxes ----
    \node[draw, dashed, rounded corners, inner sep=8pt, fit=(j1)(j3)(r)] (exactbox) {};
    \node[draw, dotted, thick, rounded corners, inner sep=12pt,
          fit=(s1)(s2)(c1)(c3)(j1)(j3)(r)] (approxbox) {};

    % ---- group labels in the side margins (no arrow/box overlap) ----
    \node[grp] at ([xshift=-13pt]approxbox.west) {Approximate join (steps 1--4)};
    \node[grp] at ([xshift=13pt]exactbox.east)   {Exact join (steps 3--4)};
  \end{tikzpicture}}
  \caption{The vector join in four steps. Steps~3--4, the in-cluster
    brute-force join and its reduction, are the \emph{exact} join; the
    \emph{approximate} join wraps them in steps~1--2, which k-means cluster the
    corpus and route each probe row to its $p$ nearest clusters so the join runs
    per cluster. The exact join of Section~\ref{sec:technical-1:-exact} is this
    inner box over a single cluster, the whole table: every probe batch is
    brute-force joined against every corpus batch, then reduced. The reduction
    (step~4) runs pushed down inside each cluster (per-row top-$k$,
    max-distance) and again globally across clusters (global top-$k$).}
  \label{fig:exact}
\end{figure}


\subsection{Super Sirius}\label{subsec:super-sirius}

Sirius is a GPU-native SQL engine.
It builds on established NVIDIA libraries such as \emph{cuDF} for columnar relational operators and \emph{RMM} for GPU memory allocation, and it ranked first on ClickBench just last year~\cite{SiriusClickBench25}.
It is designed as a drop-in accelerator, and its modular codebase is what lets new operators and features integrate cleanly.
The vector join is one such operator, and it inherits the engine's execution machinery for free.

Super Sirius is the execution layer beneath Sirius~\cite{SuperSiriusDocs}.
It compiles a query into pipelines of physical operators and runs them as tasks over GPU and CPU thread pools, and several of its components map directly onto the building blocks a vector join needs.
The \emph{task creator} turns each ready batch pair into its own task, so the join's independent pairs become independent units of work with no scheduling logic of our own.
The \emph{pipeline executor} runs those tasks concurrently across a thread pool, each on its own CUDA stream, so the GPU stays saturated with pair computations; it also reserves memory per task, so a pair that does not fit waits or spills while the query keeps running.
The \emph{task scheduler} dispatches the tasks over the available threads, and on a multi-GPU node it spreads them across GPUs in round-robin, so the pair grid fans out over the whole machine.
The \emph{downgrade executor} owns the lifetime of the batches flowing between operators and spills them across GPU, host, and disk tiers under memory pressure by way of cuCascade, so the join stays within a bounded footprint.
Finally, the \emph{result collector} streams the result back to clients in batches, so a larger-than-memory result never has to materialize at once.

\subsection{Design Overview}\label{subsec:exact-design}

We began with the exact join.
It is the most straightforward to build, and it is the foundation the approximate join extends in the next section.
For now the join is reached through a table function and requires both tables pinned in GPU memory.
This is an implementation convenience we return to in Section~\ref{subsec:exact-native}: once a vector join gains a matchable SQL surface, the engine feeds and tiers its batches the same way it does for every other operator.

The join returns one of the three reduction outcomes from Section~\ref{subsec:the-vector-join}, the $k$ nearest per probe row, the $k$ nearest pairs overall, or every pair within a threshold $\varepsilon$, and we express all three as Sirius operators.
When the corpus fits in GPU memory, the exact join is a brute-force walk over batches (Figure~\ref{fig:exact}).
Sirius reads both tables as batches, so the probe side arrives as a batch of query rows rather than a single query vector, and a brute-force search over one such batch is already a top-$k$ join between it and one corpus batch.
The join visits every pair of a probe batch and a corpus batch, runs top-$k$ on the pair, and merges each probe row's results across all corpus batches into its final top-$k$.

Our work centers on \texttt{cuvs::neighbors::brute\_force::search}, which has three properties that make it the right primitive.
The first is that it is inherently batched, which fits the vector join directly.
The second is internal tiling: the search never forms the full pairwise distance matrix between the two batches; it splits the work into tiles, computes one tile of distances at a time, selects that tile's top-$k$, and merges across tiles.
This keeps memory bounded, since only a single tile is resident at any moment rather than the whole probe-by-corpus block, so the intermediate cannot blow up as the corpus grows.
And the tile size is chosen to fit the GPU's resources, so each tile runs at high occupancy and the search stays fast.
The third property is that the distance computation maps onto GEMM: for the common metrics a tile of pairwise distances reduces to a dense matrix multiply of the probe and corpus vectors, followed by a norm correction, which is the operation the GPU runs fastest.
The one limitation is that the API serves only a top-$k$ search, whose output size is fixed per query.
The threshold join has no such bound, so no top-$k$ API fits it.
We therefore wrote our own kernel that mirrors the same tiled, GEMM-based structure and swaps the per-tile top-$k$ selection for a threshold filter.

\subsection{Operators and Search Modes}\label{subsec:exact-search-modes}

A relational join has three moving parts: a predicate that ranks the candidate pairs, a reduction that keeps the survivors, and a projection that fetches their columns.
The vector join has the same anatomy, and we give each part its own operator.
The \emph{select} operator is the leaf and carries the predicate: it holds the batch-pair grid and runs the per-pair brute-force search that ranks each pair by distance, so the heavy GEMM-bound work stays parallel and confined to the leaf.
The \emph{reduce} operator performs the per-row reduction and merges a probe batch's per-pair partials into each row's surviving top-$k$.
For global mode the reduction happens in a downstream pipeline, using Sirius's existing \texttt{TOP\_N} and \texttt{MERGE\_TOP\_N} operators.
The \emph{materialize} operator is the projection and runs last, so only the pairs that survive the reduction pull their output columns, the same discipline as pushing a projection below a join.

\paragraph{Per-row} Per-row returns the $k$ nearest corpus rows for each probe row, and it is the base the other two modes build on.
In the select operator, each batch pair goes through cuVS: \texttt{brute\_force::build} takes the corpus batch as a lightweight index and precomputes its vector norms so the search need not recompute them, and \texttt{brute\_force::search} returns, for each probe row in the pair, its $k$ nearest rows within that corpus batch.
This search is where the GEMM and the internal tiling live, so a pair is scored without ever forming its distance block in full.
Each probe row then holds one top-$k$ per corpus batch, and in the reduce operator \texttt{knn\_merge\_parts} merges those per-batch lists into the row's final top-$k$ over the whole corpus.

\paragraph{Global} Global returns the $k$ nearest pairs across the whole join.
It builds on the same per-pair search, but the reduce operator forwards the per-pair candidates without a per-row merge, and the engine's existing \texttt{TOP\_N} and \texttt{MERGE\_TOP\_N} operators reduce them into the single global top-$k$.
The materialize operator runs after this reduction and gathers columns for only the $k$ surviving pairs.

\paragraph{Threshold} Threshold returns every pair whose distance falls within a radius $\varepsilon$.
Its output size cannot be bounded ahead of time as top-$k$ does, and the number of pairs within $\varepsilon$ depends on the data and the radius.
cuVS's search APIs assume a bounded output, so none of them fits a radius join directly.
We therefore wrote our own kernel that mirrors \texttt{brute\_force::search} and replaces the top-$k$ selection with a threshold filter (\texttt{thrust::copy\_if}) that keeps every pair whose distance is within $\varepsilon$.
This preserves the two properties we want, internal tiling and GEMM, while removing the fixed output count that a top-$k$ search imposes.
The result is an edge list, one row per surviving (probe, corpus) pair carrying the two row ids and their distance.
The output buffer starts empty and grows geometrically as each tile's survivors are appended, which works the same way as a dynamic array, so the total copy work stays linear in the number of surviving pairs and no counting pass is needed.
Each tile is filtered the moment it is computed, so only one tile and the surviving pairs sit in memory at any point.
Cosine runs through the same GEMM path with no need to normalize the vectors first.

\subsection{Plain SQL and Discussion}\label{subsec:exact-native}

\begin{figure*}[t]
\centering
\tikzset{
  op/.style={draw, rounded corners, align=center, font=\scriptsize\sffamily,
             minimum width=2.1cm, minimum height=0.5cm, fill=white},
  top/.style={draw, rounded corners, align=center, font=\scriptsize\sffamily,
              minimum width=4.6cm, fill=black!5},
  flow/.style={-{Latex[length=1.5mm]}},
  blabel/.style={font=\scriptsize\bfseries\sffamily, inner sep=1pt}}

% ---- (a) distance-predicate join ----
\begin{minipage}[t]{0.48\textwidth}
\centering
\begin{lstlisting}[language=SQL, basicstyle=\ttfamily\scriptsize, frame=none,
                   aboveskip=0pt, belowskip=4pt]
SELECT *, array_distance(l.v, r.v) AS dist
FROM l JOIN r ON array_distance(l.v, r.v) <= 5;
\end{lstlisting}
\begin{tikzpicture}
  \node[op] (s0) at (0,0)     {GPU\_SCAN};
  \node[op] (s3) at (2.7,0)   {GPU\_SCAN};
  \node[op] (p1) at (0,1.3)   {PARTITION};
  \node[op] (p4) at (2.7,1.3) {PARTITION};
  \node[op] (c2) at (0,2.6)   {CONCAT};
  \node[op] (c5) at (2.7,2.6) {CONCAT};
  \node[top] (t) at (1.35,4.3)
    {RESULT\_COLLECTOR\\ PROJECTION\\ \textbf{VECTOR\_THRESHOLD\_JOIN}\\[1pt]
     {\scriptsize type: INNER}};
  \draw[flow] (s0) -- (p1);
  \draw[flow] (s3) -- (p4);
  \draw[flow] (p1) -- (c2);
  \draw[flow] (p4) -- (c5);
  \draw[flow] (c2) -- (t.south -| c2) node[blabel, midway, left]  {FULL};
  \draw[flow] (c5) -- (t.south -| c5) node[blabel, midway, right] {FULL};
\end{tikzpicture}

{\small (a) distance-predicate join}
\end{minipage}
\hfill
% ---- (b) equi-join ----
\begin{minipage}[t]{0.48\textwidth}
\centering
\begin{lstlisting}[language=SQL, basicstyle=\ttfamily\scriptsize, frame=none,
                   aboveskip=0pt, belowskip=4pt]
SELECT l.id, r.id
FROM l JOIN r ON l.id = r.id;
\end{lstlisting}
\begin{tikzpicture}
  \node[op] (s0) at (0,0)     {GPU\_SCAN};
  \node[op] (s3) at (2.7,0)   {GPU\_SCAN};
  \node[op] (p1) at (0,1.3)   {PARTITION};
  \node[op] (p4) at (2.7,1.3) {PARTITION};
  \node[op] (c2) at (0,2.6)   {CONCAT};
  \node[op] (c5) at (2.7,2.6) {CONCAT};
  \node[top] (t) at (1.35,4.3)
    {RESULT\_COLLECTOR\\ PROJECTION\\ \textbf{HASH\_JOIN}\\[1pt]
     {\scriptsize type: INNER}};
  \draw[flow] (s0) -- (p1);
  \draw[flow] (s3) -- (p4);
  \draw[flow] (p1) -- (c2);
  \draw[flow] (p4) -- (c5);
  \draw[flow] (c2) -- (t.south -| c2) node[blabel, midway, left]  {PARTIAL};
  \draw[flow] (c5) -- (t.south -| c5) node[blabel, midway, right] {PARTIAL};
\end{tikzpicture}

{\small (b) equi-join}
\end{minipage}
\caption{The same query shape under two join predicates. The distance-predicate
join~(a) and the equi-join~(b) produce identical operator trees; only the join
node and its input barriers differ. The vector threshold join consumes both
sides in full because it walks every batch pair, whereas the hash join streams
its probe side with a partial barrier.}
\label{fig:native-plans}
\end{figure*}


% again supporting our argument that vector join is a relational join
Currently only the threshold join reaches the engine through plain SQL with no special syntax.
As the Figure~\ref{fig:native-plans} shows, its operator tree is identical to a hash join on the same query: a vector join behaves like an ordinary relational join, running the same plan with a different kernel at the join node.
The one plan-level difference is the barrier, as the threshold join consumes both sides in full since it walks every batch pair, whereas the hash join streams its probe side.

% discussion about what if the tables don't fit in gpu memory
Because it has a SQL surface, the join arrives through the engine's ordinary path: its inputs are fed by the scanner as batches and tiered by the downgrade executor.
Neither table has to be pinned or fit in GPU memory, because the engine streams and spills the batches as it does for every other operator.

% lastly our situation now
Top-$k$ per-row and global do not have this surface yet.
Their syntax sugar lowers to a lateral join, whose plan shape the interceptor cannot match, so today they are reached through a table function that requires both tables to be pinned in GPU memory.
